
\documentclass[a4paper,10pt]{report}

%\usepackage[landscape]{geometry}
\usepackage[cm-default]{fontspec}
\setromanfont{FreeSerif}
\setsansfont{FreeSans}
\setmonofont{FreeMono}
\usepackage[utf8x]{inputenc}
\usepackage{fontspec}
\usepackage{xunicode}
\usepackage{xltxtra}
\usepackage{xgreek}
\usepackage{amsmath}
\usepackage{unicode-math}
\usepackage{ulem}
\usepackage{color}
\usepackage{verbatim}
\usepackage{nopageno}
\usepackage{graphicx}
\usepackage{textpos}
\setlength{\TPHorizModule}{1cm}
\setlength{\TPVertModule}{1cm}
%\usepackage[colorgrid,texcoord]{eso-pic}
\usepackage[outline]{contour}
\usepackage{wrapfig}
\usepackage{url}
\usepackage{color}
\usepackage{tikz}
\usepackage{fancybox,fancyhdr}
\usepackage{subfigure}
\usepackage{pstricks}
\usepackage{epsfig}
\usepackage{multicol}
\usepackage{listings}
\usepackage{enumerate}
\usepackage{hyperref}
\hypersetup{
  bookmarks=true,
  bookmarksopen=true,
  pdfborder=false,
  pdfpagemode=UseNone,
  raiselinks=true,
  pdfhighlight={/P},
  colorlinks,
  citecolor=black,
  filecolor=black,
  linkcolor=black,
  urlcolor=black
}


\usepackage{fontspec}
\usepackage{xunicode}
\usepackage{xltxtra}


% Margins
%\setlength{\textwidth}{24cm}

%\setlength{\voffset}{0in}
%\setlength{\textheight}{6.8in}
% Colors
\definecolor{cbrown}{rgb}{0.49,0.24,0.07}
\definecolor{cmpez}{rgb}{0.92,0.88,0.79}
\definecolor{cmpez2}{rgb}{0.62,0.33,0.34}
\definecolor{cwhite}{rgb}{1,1,1}
\definecolor{cblack}{rgb}{0,0,0}
\definecolor{cred}{rgb}{0.9,0.15,0.15}
\definecolor{clblue}{rgb}{0.57,0.85,0.97}
\definecolor{corange}{rgb}{0.86,0.49,0.18}
\definecolor{clorange}{rgb}{0.95,0.84,0.61}
\definecolor{cyellow}{rgb}{1,0.95,0.16}
\definecolor{cgreen}{rgb}{0.69,0.8,0.31}
\definecolor{clbrown}{rgb}{0.78,0.67,0.43}
\definecolor{cmagenta}{rgb}{0.79,0.05,0.54}
\definecolor{cgray}{rgb}{0.76,0.73,0.63}
\definecolor{bluesite}{rgb}{0.05,0.43,0.67}
\definecolor{redsite}{rgb}{0.90,0.15,0.15}
\definecolor{backsite}{rgb}{0.07,0.07,0.07}
% ----------- from Costas -----------------------------

%---------------   Main Fonts    -----------------%   

% \setmainfont[Mapping=tex-text]{Calibri}
% \setmainfont[Mapping=tex-text]{Times New Roman}
%\setmainfont[Mapping=tex-text]{Droid Serif}
%\setmainfont[Mapping=tex-text]{Cambria}
%\setmainfont[Mapping=tex-text]{cm-unicode}
% \setmainfont[Mapping=tex-text]{Gentium}
% \setmainfont[Mapping=tex-text]{GFS Didot}
% \setmainfont[Mapping=tex-text]{Comic Sans MS}
 \setmainfont[Mapping=tex-text]{Ubuntu}
% \setmainfont[Mapping=tex-text]{Myriad Pro}
% \setmainfont[Mapping=tex-text]{CMU Concrete}
%\setmainfont[Mapping=tex-text]{DejaVu Sans}
%\setmainfont[Mapping=tex-text]{KerkisSans}
%\setmainfont[Mapping=tex-text]{KerkisCaligraphic}
% \setmainfont[Mapping=tex-text]{Segoe Print}
% \setmainfont[Mapping=tex-text]{Gabriola}

%------------------------------------------------%

% ------------ Mathematics Fonts ----------------%
\setmathfont{Asana-Math.ttf}
% -----------------------------------------------%
% Misc
\author{Κασωτάκη Ε. - Σμαραγδάκης Κ.}
\title{www.math24.gr}
% ---------------- Formation --------------------%


\setlength\topmargin{-0.7cm}
\addtolength\topmargin{-\headheight}
\addtolength\topmargin{-\headsep}
\setlength\textheight{26cm}
\setlength\oddsidemargin{-0.54cm}
\setlength\evensidemargin{-0.54cm}
%\setlength\marginparwidth{1.5in}
\setlength\textwidth{17cm}

\RequirePackage[avantgarde]{quotchap}
\renewcommand\chapterheadstartvskip{\vspace*{0\baselineskip}}
\RequirePackage[calcwidth]{titlesec}
\titleformat{\section}[hang]{\bfseries}
{\Large\thesection}{12pt}{\Large}[{\titlerule[0.9pt]}]
%--------------------------------------------------%
\usepackage{draftwatermark}
\SetWatermarkText{www.math24.gr}
\SetWatermarkLightness{0.9}
\SetWatermarkScale{0.6}
%--------------------------------------------------%


\usepackage{xcolor}
\usepackage{amsthm}
\usepackage{framed}
\usepackage{parskip}

\colorlet{shadecolor}{bluesite!20}

\newtheorem*{orismos}{Ορισμός}

\newenvironment{notation}
  {\begin{shaded}\begin{theorem}}
  {\end{theorem}\end{shaded}}




% Document begins
\begin{document}
\pagestyle{fancy}
\fancyhead{}
\fancyfoot{}
\renewcommand{\headrulewidth}{0pt}
\renewcommand{\footrulewidth}{0pt}


\fancyhead[LO,LE]{
 %\textblockcolor{backsite}
 %\begin{textblock}{5}(-2,-0.55)
  %\rule{0cm}{1cm}
 %\end{textblock}
 \textblockcolor{bluesite}
 \begin{textblock}{5}(-1.5,-0.55)
  \rule{0cm}{1cm}
 \end{textblock}
 %\textblockcolor{bluesite}
 %\begin{textblock}{14}(5.5,-0.55)
 % \rule{0cm}{1cm}
 %\end{textblock}
 \begin{textblock}{0}(-1,-0.25)
 \color{cwhite} \begin{Large}www.math24.gr\end{Large}
 \end{textblock}
\begin{textblock}{0}(16,-0.4)
 \color{cwhite} \includegraphics[height=1.5cm]{math24_logo.png}
 \end{textblock}
\textblockcolor{white}
\begin{textblock}{17}(-1,28)
 \color{backsite} \begin{small}Copyright \textcopyright 2013, Κασωτάκη Ε.(ikasotaki@gmail.com) - Σμαραγδάκης Κ.(kesmarag@gmail.com)\end{small}
 \end{textblock}
}
 
\begin{shaded}
\begin{center}
\huge \textbf{Φυλλάδιο Ασκήσεων}\\
%Πρόσθεση ρητών αριθμών
\end{center} 
\textbf{Μαθηματικά Α' Γυμνασίου} \hfill \textbf{Ημερομηνία Παράδοσης : \hspace{2em} }
\subsection*{Ονοματεπώνυμο :\hfill  \hspace{5em}}
\end{shaded}
\vspace{2em}
\begin{itemize}
 \item Θετικοί αριθμοί
 \item Αρνητικοί αριθμοί
 \item Εφαρμογή θετικών και αρνητικών αριθμών στην καθημερινή ζωή
 \item Ομόσημοι και Ετερόσημοι αριθμοί
 \item Ακέραιοι και Ρητοί αριθμοί
\end{itemize}
\section*{Θεωρία - Θετικοί και Αρνητικοί αριθμοί\hfill \small{}}
\begin{itemize}
 \item \textbf{Πρόσημα} ονομάζονται τα σύμβολα \textbf{"+"} και \textbf{"-"}.              
 \item \textbf{Θετικοί αριθμοί} ονομάζονται οι αριθμοί που μπροστά τους έχουν το πρόσημο \textbf{"+"}. \\
        \textbf{π.χ} ο αριθμός $+3 $ \\
       \textbf{π.χ} ο αριθμός $+7 $\\
       \textbf{π.χ} ο αριθμός $+13 $
 \item \textbf{Αρνητικοί αριθμοί} ονομάζονται οι αριθμοί που μπροστά τους έχουν το πρόσημο \textbf{"-"}.\\
        \textbf{π.χ} ο αριθμός $-1 $ \\
       \textbf{π.χ} ο αριθμός $-\dfrac{5}{2} $\\
       \textbf{π.χ} ο αριθμός $-47 $
\end{itemize}
Σημείωση:
\begin{itemize}
 \item Αν ένας αριθμός δεν έχει μπροστά του κανένα πρόσημο, τότε εννοείται το πρόσημο \textbf{"+"}, δηλαδή 
       ο αριθμός είναι θετικός.\\
       \textbf{π.χ} $41$ είναι το ίδιο με το $+41$ 
 \item Το μηδέν δεν είναι ούτε θετικός ούτε αρνητικός αριθμός.
\end{itemize}



\section*{Άσκηση 1  (Αντίθετοι αριθμοί)\hfill \small{25 μονάδες}}
Να κατατάξετε τους παρακάτω αριθμούς σε δύο ομάδες, τους θετικούς και τους αρνητικούς:
$$ -3 \quad ,\quad  5,1 \quad,\quad -7\quad,\quad -\dfrac{1}{2}\quad,\quad +1\quad,\quad +\dfrac{11}{3}\quad,\quad 7    $$


\section*{Εφαρμογή θετικών και αρνητικών αριθμών στην καθημερινή ζωή\hfill \small{}}
\begin{itemize}
 \item  \textbf{Οι θετικοί αριθμοί} χρησιμοποιούνται για να εκφράσουν
        \begin{itemize}
         \item πάνω (από κάτι) \textbf{π.χ} 10\textdegree πάνω από το 0
         \item άνοδος, αύξηση \textbf{π.χ} ο μισθός αυξήθηκε κατά 53\texteuro
         \item κέρδος \textbf{π.χ} ο μπακάλης είχε μηνιαίο κέρδος 720 \texteuro
        \end{itemize}
 \item  \textbf{Οι αρνητικοί αριθμοί} χρησιμοποιούνται για να εκφράσουν
        \begin{itemize}
         \item κάτω (από κάτι) \textbf{π.χ} 2 βαθμοί Κελσίου κάτω από το 0: -2\textdegree C
         \item πτώση, μείωση, ζημιά, έκπτωση \textbf{π.χ} έκπτωση 5\texteuro : -5\texteuro
        \end{itemize}
\end{itemize}

\section*{Άσκηση 2  \hfill \small{25 μονάδες}}
Να εκφράσετε με τη βοήθεια των θετικών και αρνητικών αριθμών τις παρακάτω προτάσεις :
\begin{enumerate}[1)]
\item $3$ βαθμοί Κελσίου κάτω από το μηδέν
\item $3,4m$ πάνω από την επιφάνεια της θάλασσας
\item Ο Θανάσης ψήλωσε κατά $2cm$
\item Ο μισθός αυξήθηκε κατά $105$\texteuro
\item Τα παπούτσια είχαν έκπτωση $19$\texteuro
\end{enumerate}


\section*{Θεωρία - Ομόσημοι και Ετερόσημοι αριθμοί\hfill \small{}}
\begin{itemize}
 \item \textbf{Ομόσημοι} λέγονται οι αριθμοί που έχουν το ίδιο πρόσημο\\
       \textbf{π.χ} $+3$ και $+7,1$\\
       \textbf{π.χ} $-1$ και $-7$\\
       \textbf{π.χ} $13$, $110$ και $241$
 \item \textbf{Ετερόσημοι} λέγονται οι αριθμοί που έχουν το διαφορετικό πρόσημο\\
       \textbf{π.χ} $+17$ και $-23$\\
       \textbf{π.χ} $-1$ και $+71$
\end{itemize}


\section*{Άσκηση 3  \hfill \small{25 μονάδες}}
Να συμπληρώσετε τον παρακάτω πίνακα σημειώνοντας \textbf{"ΝΑΙ"} ή \textbf{"ΟΧΙ"} σε κάθε κενό:
\begin{center}
\begin{tabular}{|c|c|c|}
\hline
\textbf{Αριθμοί}  & \textbf{Ομόσημοι} & \textbf{Ετερόσημοι}\\ \hline
$43$ και $+7$ &  & \\ \hline
$-2$ και $+2$ &  & \\ \hline
$13$ και $-5$ &  & \\ \hline 
$-3$,$-5$ και $-9$ &  & \\ \hline
$1$, $11$ και $21$ &  & \\ \hline
\end{tabular}
\end{center}


\section*{Θεωρία - Ακέραιοι και Ρητοί αριθμοί\hfill \small{}}
\begin{itemize}
 \item \textbf{Φυσικοί αριθμοί} είναι οι αριθμοί $0,1,2,3,4,\cdots ,11,12,\cdots , 997, \cdots$ \\
       Κάθε φυσικός αριθμός έχει έναν επόμενο και έναν προηγούμενο αριθμό (εκτός από το $0$ που έχει 
       μόνο επόμενο, το $1$)
 \item \textbf{Ακέραιοι αριθμοί} είναι οι φυσικοί αριθμοί μαζί με τους αντίστοιχους αρνητικούς αριθμούς.\\
        Δηλαδή $\cdots ,-4,-3,-2,-1,0,1,2,3,4,\cdots$
\item \textbf{Ρητοί αριθμοί} είναι οι φυσικοί αριθμοί, τα κλάσματα και οι δεκαδικοί αριθμοί μαζί με τους 
      αντίστοιχους αρνητικούς αριθμούς \\
      \textbf{π.χ} $1,+7,-2, \dfrac{5}{3}, +532$
\end{itemize}




\section*{Άσκηση 4  \hfill \small{25 μονάδες}}
Να συμπληρώσετε τον παρακάτω πίνακα σημειώνοντας \textbf{"ΝΑΙ"} ή \textbf{"ΟΧΙ"} σε κάθε κενό:
\begin{center}
\begin{tabular}{|c|c|c|c|}
\hline
\textbf{Αριθμός}  & \textbf{Φυσικός} & \textbf{Ακέραιος} & \textbf{Ρητός}\\ \hline
$1$ &  & &\\ \hline
$-2$ &  & &\\ \hline
$-2,3$  &  & &\\ \hline 
$-\dfrac{5}{2}$ & &  & \\ \hline
$3$ & &  & \\ \hline
$13,9$ &  & &\\ \hline
$-41$ &  & &\\ \hline
$+317$  &  & &\\ \hline 
$+\dfrac{4}{5}$ & &  & \\ \hline
$1,2$ & &  & \\ \hline
\end{tabular}
\end{center}







\end{document}
